É com profunda gratidão que dedico minha Monografia às pessoas que estiveram ao meu lado, especialmente minha família, que foi a base para que eu chegasse até esse ponto na minha graduação, e ao meu orientador, que desde o início do meu curso me inspirou a ser uma pessoa comprometida com a qualidade, me preparando para um mundo que não aceita nada pela metade, mas espera sempre pelo ``algo mais'' que faz a verdadeira diferença.

Agradeço às pessoas que não estiveram ao meu lado, em especial Aaron Swartz e Alexandra Elbakyan, por serem símbolos do movimento do acesso aberto à produção científica, iniciativa que hoje deu frutos ao me possibilitar acesso a diversos artigos e subsidiou teoricamente todo esse trabalho. A Taiichi Ohno e Kiichiro Toyoda, por influenciarem dramaticamente a minha ética de trabalho em tudo que me disponho a fazer com a filosofia que hoje conhecemos como \textit{Lean Manufacturing}.

Agradeço também a todos que passaram por minha vida de forma passageira e deixaram grandes marcas. Em especial, à enfermeira que cuidou de mim no Hospital Santa Izabel e me mencionou a importância da gestão e de comunicação no âmbito profissional. Essa filosofia me inspirou nesse trabalho como um todo, a partir do momento que me comprometi com o planejamento e com o feedback de todas as partes interessadas no decorrer desse trabalho.

Esta monografia não é apenas um marco acadêmico, mas uma realização coletiva. A todos vocês, o meu mais sincero agradecimento.
